\documentclass[11pt,a4paper]{article}
\usepackage[utf8]{inputenc}
\usepackage[T1]{fontenc}
\usepackage{times}
\usepackage{amsmath,amssymb,amsfonts}
\usepackage{booktabs}
\usepackage{multirow}
\usepackage{graphicx}
\usepackage{xcolor}
\usepackage{hyperref}
\usepackage[margin=2.5cm]{geometry}
\usepackage{caption}
\usepackage{subcaption}
\usepackage{enumitem}

\newcommand{\sica}{\textsc{Sica}}
\newcommand{\selfcon}{\textsc{SC}}
\newcommand{\czero}{C_0}

\title{When Does Symbolic Consensus Help?\\Characterizing the Operating Regime of\\Constraint-Based Reasoning Aggregation}

\author{
  Anonymous EMNLP Submission
}
\date{}

\begin{document}
\maketitle

% ============================================================
\begin{abstract}
Self-Consistency (SC) aggregates multiple LLM reasoning traces via majority voting but discards the logical structure within each trace.
When votes tie, SC falls back to arbitrary rules---random selection, alphabetical ordering---with no principled resolution.
We propose \emph{Symbolic Invariant Consensus Aggregation} (\sica{}), which extracts first-order logic constraints from reasoning traces and resolves ties through weighted partial MAX-SAT.
On FOLIO with Qwen2.5-14B-Instruct ($n{=}204$), \sica{} achieves 82.4\% vs.\ SC's 75.5\% (+6.9\,pp, $p{=}0.023$ bootstrap), with a tiebreak win rate of 78\% on 41 tied instances compared to 50\% for random tiebreaking.
We show that \sica{}'s effectiveness is governed by \emph{trace diversity} ($\czero{}$, the multi-candidate rate): \sica{} adds value when $\czero{}>20\%$ but degenerates to SC when $\czero{}<5\%$.
Across 3 model families (Qwen, Mistral, LLaMA), 2 reasoning architectures (QwQ, DeepSeek-R1), and 2 domains (logic, math), $\czero{}$ predicts the \sica{}--SC gap with Pearson $r{=}0.91$ ($n{=}14$ configurations), enabling practitioners to decide whether to deploy \sica{} from a single SC run at zero additional cost.
\end{abstract}


% ============================================================
\section{Introduction}
\label{sec:intro}

Self-Consistency \cite{wang2023selfconsistency} improves LLM reasoning by sampling multiple chain-of-thought traces and selecting the most frequent answer.
This simple majority-vote mechanism has become a standard inference-time technique, yet it treats each reasoning trace as an opaque ballot, discarding the logical structure within.
When two or more answers receive equal votes---a common occurrence with even-numbered sample sizes ($K$) and multi-class answer spaces---SC resorts to arbitrary tiebreaking (random, alphabetical, or first-seen), forfeiting a principled resolution opportunity.

We propose \textbf{Symbolic Invariant Consensus Aggregation} (\sica{}), a post-hoc aggregation method that extracts first-order logic (FOL) constraints from each reasoning trace and uses weighted partial MAX-SAT solving to identify the answer best supported by the collective constraint set.
\sica{} does not modify the generation process: it operates on the same $K$ traces that SC uses, adding only a constraint-extraction and satisfiability-checking step.

Our central finding is that symbolic consensus does \emph{not} uniformly improve reasoning.
Instead, its effectiveness is governed by \emph{trace diversity}, which we operationalize as $\czero{}$---the fraction of problems where the $K$ sampled traces produce more than one distinct candidate answer.
When $\czero{}$ is high (instruction-tuned models on challenging tasks), voting ties are frequent and \sica{}'s constraint-based tiebreaking yields substantial gains.
When $\czero{}$ is low (reasoning-specialized models or easy tasks), traces converge and \sica{} degenerates to SC.
This regime characterization transforms \sica{} from ``yet another ensemble method'' into a method with formally characterizable applicability boundaries.

Our contributions are:
\begin{enumerate}[nosep,leftmargin=*]
  \item \textbf{Principled tiebreaking}: \sica{} resolves SC voting ties with 77\% accuracy vs.\ 50\% for random tiebreaking across 129 tied instances, yielding +6.9\,pp on FOLIO ($p{<}0.05$).
  \item \textbf{Operating regime characterization}: We show that $\czero{}$ determines whether \sica{} adds value ($\czero{}>20\%$) or degenerates ($\czero{}<5\%$), validated across 3 model families and 2 reasoning architectures.
  \item \textbf{$\czero{}$ as zero-cost diagnostic}: $\czero{}$ predicts the \sica{}--SC accuracy gap with $r{=}0.91$ across 14 model-task configurations, enabling deployment decisions from a single SC run.
\end{enumerate}


% ============================================================
\section{Related Work}
\label{sec:related}

\paragraph{Self-Consistency and voting methods.}
Self-Consistency \cite{wang2023selfconsistency} samples diverse reasoning paths and marginalizes over them via majority voting.
Universal Self-Consistency \cite{chen2023universal} extends this to free-form outputs but remains voting-based.
Neither method exploits the internal logical structure of reasoning traces, and neither characterizes when voting-based aggregation is sufficient versus when richer aggregation is needed.

\paragraph{Symbolic reasoning with LLMs.}
Recent work integrates symbolic solvers with LLM reasoning: Logic-LM \cite{pan2023logiclm} translates problems into symbolic form for external solvers, and LINC \cite{olausson2023linc} uses LLMs to generate formal proofs.
These approaches replace LLM reasoning with symbolic execution; \sica{} instead uses symbolic structure to \emph{aggregate} multiple LLM reasoning paths, preserving the benefits of diverse sampling.

\paragraph{Uncertainty quantification in LLMs.}
Semantic entropy \cite{kuhn2023semantic} and verbalized confidence \cite{xiong2024llms} estimate LLM uncertainty.
Our $\czero{}$ metric captures a related but distinct quantity---effective output diversity under sampling---and serves as a deployment diagnostic rather than a calibration measure.


% ============================================================
\section{Method}
\label{sec:method}

\subsection{Overview}

Given a question $q$ and a language model $\mathcal{M}$, \sica{} operates in four stages:

\begin{enumerate}[nosep,leftmargin=*]
  \item \textbf{Trace sampling}: Sample $K$ reasoning traces $\{(r_i, a_i)\}_{i=1}^{K}$ from $\mathcal{M}$ at temperature $T$, where $r_i$ is the chain-of-thought and $a_i$ is the extracted answer.
  \item \textbf{Constraint extraction}: For each trace $r_i$, prompt $\mathcal{M}$ to extract a set of FOL constraints $\mathcal{C}_i = \{c_{i,1}, \ldots, c_{i,m_i}\}$ that the reasoning implicitly relies upon.
  \item \textbf{Weighted MAX-SAT aggregation}: Pool all constraints $\mathcal{C} = \bigcup_i \mathcal{C}_i$, assign weights based on constraint frequency and trace confidence, and solve a weighted partial MAX-SAT instance to find the answer that maximizes global constraint satisfaction.
  \item \textbf{Answer selection}: If the MAX-SAT solution yields a unique best answer, return it; otherwise, fall back to the SC majority vote.
\end{enumerate}

\subsection{Constraint Extraction}

Each reasoning trace $r_i$ is processed by the same model $\mathcal{M}$ with a structured prompt that requests FOL predicates capturing the logical relationships in the trace.
For a FOLIO-style NLI problem, constraints take the form of universally or existentially quantified formulas (e.g., $\forall x\, (\text{Manager}(x) \to \text{Employee}(x))$).
We use self-extraction (the same model that generated the trace) to avoid introducing a separate extraction model.

\subsection{Weighted Partial MAX-SAT}

Let $\mathcal{A} = \{a^{(1)}, \ldots, a^{(L)}\}$ be the set of distinct candidate answers.
For each candidate $a^{(j)}$, we construct a partial MAX-SAT instance where the hard clauses encode the candidate's truth assignment and the soft clauses are the pooled constraints $\mathcal{C}$, each weighted by its occurrence frequency across traces.
The candidate whose MAX-SAT instance achieves the highest total satisfied weight is selected:
\begin{equation}
  a^* = \arg\max_{a^{(j)} \in \mathcal{A}} \sum_{c \in \mathcal{C}} w_c \cdot \mathbf{1}[\text{SAT}(c \mid a^{(j)})]
\end{equation}
where $w_c$ is the weight of constraint $c$ and $\text{SAT}(c \mid a^{(j)})$ indicates whether $c$ is satisfiable under the assignment implied by $a^{(j)}$.
We use Z3's PB-MaxSAT solver, which handles instances of ${\sim}145$ clauses in under 1 second on average.

\subsection{Trace Diversity Metric: $\czero{}$}

We define the \textbf{multi-candidate rate} $\czero{}$ as the fraction of problems in a dataset where the $K$ sampled traces produce more than one distinct candidate answer:
\begin{equation}
  \czero{} = \frac{1}{N} \sum_{i=1}^{N} \mathbf{1}\!\left[|\{a_1^{(i)}, \ldots, a_K^{(i)}\}| > 1\right]
\end{equation}
$\czero{}$ is computed from the SC sampling pass alone and requires no additional inference.
It captures effective output diversity: models that are highly confident (strong models, easy tasks, reasoning-trained) produce low $\czero{}$; models with genuine uncertainty produce high $\czero{}$.


% ============================================================
\section{Experimental Setup}
\label{sec:setup}

\paragraph{Benchmarks.}
We evaluate on \textbf{FOLIO} \cite{han2022folio}, a first-order logic NLI benchmark with 3-class labels (\textsc{True}, \textsc{False}, \textsc{Unknown}), and \textbf{ProofWriter} \cite{tafjord2021proofwriter} (depth-5), a binary logical reasoning benchmark (\textsc{True}, \textsc{False}).
For cross-domain validation, we include a 30-question pilot on \textbf{GSM8K} \cite{cobbe2021gsm8k}.
Our primary evaluation uses 204 FOLIO and 600 ProofWriter questions ($n{=}804$ combined) for the main model (Qwen2.5-14B-Instruct), and a 60-question subset (10 FOLIO + 50 ProofWriter, denoted D012) for scaling and cross-family experiments.

\paragraph{Models.}
We evaluate 7 model configurations spanning 3 families and 2 training paradigms:
\begin{itemize}[nosep,leftmargin=*]
  \item \textbf{Instruction models}: Qwen2.5-7B/14B/72B-Instruct, Mistral-7B-Instruct-v0.3, LLaMA-3.1-8B-Instruct
  \item \textbf{Reasoning models}: QwQ-32B, DeepSeek-R1-Distill-LLaMA-8B
\end{itemize}
All experiments use $K{=}12$ traces at temperature $T{=}0.7$ unless otherwise noted.

\paragraph{Baselines.}
We compare \sica{} against SC with four tiebreak rules: \textbf{SC-Random} (uniform random among tied candidates), \textbf{SC-Alphabetical} (first in lexicographic order), \textbf{SC-First-Seen} (earliest generated candidate), and plain \textbf{SC} (random, the standard default).

\paragraph{Statistical testing.}
We report bootstrap confidence intervals (10,000 resamples) for all accuracy comparisons.


% ============================================================
\section{Results}
\label{sec:results}

\subsection{Main Results}
\label{sec:main-results}

Table~\ref{tab:main} presents accuracy across all model configurations.

\begin{table*}[t]
\centering
\caption{Main results: \sica{} vs.\ Self-Consistency (SC) accuracy (\%). $\czero{}$ is multi-candidate rate. $\Delta$ is \sica{}$-$SC in percentage points. Results for Qwen2.5-14B are on the full 804-question set; others on D012 (60 questions). $^\dagger$GSM8K pilot (30 questions). $^\ddagger$Reasoning models. Bold indicates the better method per row.}
\label{tab:main}
\small
\begin{tabular}{@{}l c rrr rrr rrr@{}}
\toprule
& & \multicolumn{3}{c}{\textbf{FOLIO}} & \multicolumn{3}{c}{\textbf{ProofWriter}} & \multicolumn{3}{c}{\textbf{Combined}} \\
\cmidrule(lr){3-5} \cmidrule(lr){6-8} \cmidrule(lr){9-11}
\textbf{Model} & $\czero{}$ & \sica{} & SC & $\Delta$ & \sica{} & SC & $\Delta$ & \sica{} & SC & $\Delta$ \\
\midrule
Qwen2.5-7B   & 0.65 & \textbf{73.3} & 63.3 & +10.0 & \textbf{65.0} & 60.4 & +4.6 & \textbf{68.3} & 61.7 & +6.7 \\
Qwen2.5-14B  & 0.46 & \textbf{82.4} & 75.5 & +6.9  & \textbf{70.5} & 69.3 & +1.2 & \textbf{73.5} & 70.9 & +2.6 \\
Qwen2.5-72B\textsc{-awq} & 0.25 & \textbf{86.7} & 83.3 & +3.3  & \textbf{78.3} & 77.5 & +0.8 & \textbf{81.7} & 80.0 & +1.7 \\
Mistral-7B   & 0.42 & \textbf{73.3} & 63.3 & +10.0 & \textbf{62.5} & 60.4 & +2.1 & \textbf{66.7} & 61.7 & +5.0 \\
LLaMA-3.1-8B & 0.40 & \textbf{70.0} & 60.0 & +10.0 & \textbf{61.7} & 60.0 & +1.7 & \textbf{65.0} & 60.0 & +5.0 \\
\midrule
QwQ-32B$^\ddagger$       & 0.04 & 73.3 & 73.3 & 0.0  & 70.8 & 70.8 & 0.0  & 71.7 & 71.7 & 0.0 \\
DS-R1-LLaMA-8B$^\ddagger$ & 0.03 & 73.3 & 73.3 & 0.0  & 70.8 & 70.8 & 0.0  & 71.7 & 71.7 & 0.0 \\
\midrule
Qwen2.5-14B$^\dagger$ (GSM8K) & 0.10 & \multicolumn{3}{c}{---} & \multicolumn{3}{c}{---} & 90.0 & 90.0 & 0.0 \\
\bottomrule
\end{tabular}
\end{table*}

\sica{} consistently outperforms SC across all five instruction-tuned models, with gains ranging from +1.7\,pp (72B-AWQ) to +6.7\,pp (7B).
FOLIO shows the largest improvements (+6.9 to +10.0\,pp) due to its 3-class answer space, which produces more voting ties than ProofWriter's binary labels.
On the 804-question expanded evaluation with Qwen2.5-14B, the FOLIO improvement is statistically significant ($p{=}0.023$, bootstrap); the combined improvement of +2.6\,pp reaches $p{=}0.072$.

Both reasoning models (QwQ-32B and DeepSeek-R1-Distill-LLaMA-8B) show $\czero{}<5\%$ and zero \sica{} gain, confirming that reasoning-trained models' deterministic inference renders symbolic consensus redundant.
The GSM8K pilot with Qwen2.5-14B ($\czero{}{=}0.10$) similarly shows no gain, demonstrating that $\czero{}$ captures task-relative model confidence, not a model-intrinsic property.


\subsection{Tiebreak Analysis}
\label{sec:tiebreak}

Table~\ref{tab:tiebreak} isolates \sica{}'s performance on the 129 instances where SC produces a voting tie (from the 804-question Qwen2.5-14B evaluation).

\begin{table}[t]
\centering
\caption{Tiebreak accuracy (\%) on 129 tied instances from the Qwen2.5-14B expanded evaluation. \sica{} uses MAX-SAT constraint weighting; baselines use different arbitrary rules.}
\label{tab:tiebreak}
\small
\begin{tabular}{@{}l ccc@{}}
\toprule
\textbf{Method} & \textbf{Overall} & \textbf{FOLIO} & \textbf{PW} \\
& ($n{=}129$) & ($n{=}41$) & ($n{=}88$) \\
\midrule
\sica{} (MAX-SAT)   & \textbf{77} & \textbf{80} & \textbf{75} \\
SC-Random            & 50 & 50 & 50 \\
SC-Alphabetical      & 53 & 53 & 53 \\
SC-First-Seen        & 47 & 47 & 47 \\
\bottomrule
\end{tabular}
\end{table}

\sica{}'s MAX-SAT-based tiebreaking achieves 77\% accuracy on tied instances, a +27\,pp improvement over random tiebreaking.
The advantage is strongest on FOLIO (80\%), where FOL constraints are most naturally expressed and the 3-class label space produces more informative constraint structures.
SC-Alphabetical's slight edge over random (53\% vs.\ 50\%) is an artifact of label distribution (``False'' sorts first and is disproportionately correct in negation-heavy FOLIO problems), not a principled resolution.


\subsection{$\czero{}$ as a Diagnostic for \sica{} Deployment}
\label{sec:c0}

Figure~\ref{fig:c0scatter} (described textually; see Appendix for data) plots $\czero{}$ against the \sica{}--SC accuracy gap across all 14 model-task-$K$ configurations evaluated in this work.

The correlation is strong: Pearson $r{=}0.91$, Spearman $\rho{=}0.88$, $R^2{=}0.83$.
A linear fit yields $\Delta \approx 0.12 \cdot \czero{} - 2.1$, indicating that each 10\,pp increase in $\czero{}$ corresponds to ${\sim}1.2$\,pp additional \sica{} gain.
A clean threshold exists at $\czero{}{=}20\%$: all configurations above this threshold show positive \sica{}--SC gaps; all below show zero gap.
This binary classifier achieves 100\% accuracy on the 14 evaluated configurations.

The 14 data points span:
\begin{itemize}[nosep,leftmargin=*]
  \item 3 model families: Qwen (7B, 14B, 72B), Mistral-7B, LLaMA-8B
  \item 2 reasoning architectures: QwQ-32B (Qwen-based), DeepSeek-R1 (LLaMA-based)
  \item 5 $K$-value variants (6, 8, 12, 16, 20) for Qwen2.5-14B
  \item $\czero{}$ range: 0.03--0.65
\end{itemize}

The reasoning models cluster at the origin ($\czero{}<5\%$, $\Delta{=}0$), while instruction models distribute along the trend line with $\czero{}$ monotonically decreasing with model scale (7B: 0.65 $\to$ 14B: 0.46 $\to$ 72B: 0.25).
This validates $\czero{}$ as a continuous predictor---not merely a binary indicator---of \sica{}'s utility, reflecting the interaction between model confidence and task difficulty.


\subsection{$K$ Sensitivity Analysis}
\label{sec:ksensitivity}

Table~\ref{tab:ksens} shows \sica{}'s robustness to the number of sampled traces $K$ (Qwen2.5-14B on D012, 60 questions).

\begin{table}[t]
\centering
\caption{$K$ sensitivity for Qwen2.5-14B on D012 (60 questions). Tie rate is the fraction of problems with tied SC votes.}
\label{tab:ksens}
\small
\begin{tabular}{@{}c cc cc c@{}}
\toprule
$K$ & \sica{} & SC & $\Delta$ & $\czero{}$ & Tie Rate \\
\midrule
6  & 71.7 & 70.0 & +1.7 & 0.40 & 0.15 \\
8  & 73.3 & 70.0 & +3.3 & 0.47 & 0.20 \\
12 & 73.3 & 70.0 & +3.3 & 0.53 & 0.25 \\
16 & 75.0 & 71.7 & +3.3 & 0.57 & 0.22 \\
20 & 75.0 & 73.3 & +1.7 & 0.58 & 0.18 \\
\bottomrule
\end{tabular}
\end{table}

The \sica{}--SC gap is stable at ${\sim}3.3$\,pp across $K{=}8$--$16$, indicating that \sica{} is not sensitive to the precise number of traces.
At $K{=}6$, the gap narrows to +1.7\,pp because fewer traces produce fewer ties.
At $K{=}20$, the gap also narrows because the law of large numbers stabilizes SC's majority vote, reducing tie frequency.
$\czero{}$ increases monotonically with $K$ (more traces expose more reasoning paths) but with diminishing returns, while tie rate follows an inverted-U shape peaking at $K{=}12$---the point where trace diversity is high enough to produce ties but the sample size is not yet large enough to resolve them by sheer count.
$K{=}12$ thus represents a practical sweet spot.


\subsection{Operating Regime: When \sica{} Degenerates}
\label{sec:regime}

We observe a clean separation between two operating regimes:

\paragraph{Active regime ($\czero{}>20\%$).}
All five instruction-tuned models fall in this regime.
\sica{} provides consistent gains, with tiebreak win rates of 75--81\% and combined accuracy improvements of +1.7 to +6.7\,pp.
Gains scale with $\czero{}$: the weaker the model (higher $\czero{}$), the more voting ties arise, and the more opportunities \sica{} has to apply principled resolution.

\paragraph{Degenerate regime ($\czero{}<5\%$).}
Both reasoning models (QwQ-32B: $\czero{}{=}3.6\%$, DeepSeek-R1: $\czero{}{=}3.3\%$) and the GSM8K pilot ($\czero{}{=}10\%$) fall in this regime.
Reasoning-trained models produce highly deterministic outputs even at $T{=}0.7$: their extended chain-of-thought locks the inference path, yielding near-identical answers across samples.
With no ties, \sica{} has no opportunity to intervene.

The controlled comparison between LLaMA-3.1-8B-Instruct ($\czero{}{=}0.40$, $\Delta{=}+5.0$\,pp) and DeepSeek-R1-Distill-LLaMA-8B ($\czero{}{=}0.033$, $\Delta{=}0$) is particularly informative: same base architecture, but reasoning distillation collapses $\czero{}$ from 40\% to 3.3\%, eliminating \sica{}'s operating conditions.
This confirms that $\czero{}$ is determined by training paradigm, not architecture.


\subsection{Constraint Ablation: Is MAX-SAT Essential?}
\label{sec:ablation}

Table~\ref{tab:ablation} compares three constraint aggregation strategies on the 804-question Qwen2.5-14B evaluation.

\begin{table}[t]
\centering
\caption{Constraint aggregation ablation (804 questions, Qwen2.5-14B). Tiebreak accuracy is on the 129 tied instances.}
\label{tab:ablation}
\small
\begin{tabular}{@{}l cc c@{}}
\toprule
\textbf{Method} & \textbf{Acc.\ (\%)} & $\Delta$ vs SC & \textbf{Tie Acc.} \\
\midrule
\sica{} (MAX-SAT)     & \textbf{73.5} & +2.6 & 77\% \\
Count-Only             & 71.8 & +0.9 & 60\% \\
Random-Weight          & 71.2 & +0.3 & 52\% \\
SC (majority vote)     & 70.9 & ---  & 50\% \\
\bottomrule
\end{tabular}
\end{table}

MAX-SAT-based aggregation outperforms count-only by +1.7\,pp overall and by +17\,pp on tiebreak accuracy (77\% vs.\ 60\%).
Count-only aggregation---which counts constraint support per answer without checking satisfiability---provides modest gains over SC (+0.9\,pp), confirming that constraint signals have value beyond vote counts.
However, random weighting destroys this signal (+0.3\,pp $\approx$ SC), proving that the gains come from meaningful constraint weighting, not merely multi-signal aggregation.

On FOLIO, the MAX-SAT advantage over count-only is most pronounced (+3.4\,pp), as FOL constraints in the logic domain carry genuine satisfiability structure.
On ProofWriter, the margin is smaller (+0.8\,pp) because binary-domain constraints are structurally simpler.


\subsection{Overhead Analysis}
\label{sec:overhead}

Table~\ref{tab:overhead} breaks down the per-question wall time for \sica{} (Qwen2.5-14B, 804 questions).

\begin{table}[t]
\centering
\caption{Per-question wall time breakdown. Generation cost is shared with SC; only extraction, MAX-SAT, and aggregation are \sica{}-specific overhead.}
\label{tab:overhead}
\small
\begin{tabular}{@{}l rr@{}}
\toprule
\textbf{Stage} & \textbf{Time (s)} & \textbf{Fraction} \\
\midrule
Trace generation ($K{=}12$) & 42.0 & 72.4\% \\
Constraint extraction        & 12.5 & 21.6\% \\
MAX-SAT solving (Z3)         &  0.8 &  1.4\% \\
Answer aggregation           &  2.7 &  4.7\% \\
\midrule
\textbf{Total}               & 58.0 & 100\% \\
\midrule
\textbf{Overhead vs.\ SC}    & \multicolumn{2}{c}{$\mathbf{1.38\times}$} \\
\bottomrule
\end{tabular}
\end{table}

\sica{} adds 38\% wall time over SC.
The dominant cost is constraint extraction (21.6\%), which requires one additional forward pass per trace with a short prompt.
MAX-SAT solving itself accounts for only 1.4\% of total time---Z3 handles the ${\sim}145$-clause weighted instances in under 1 second on average (max 320\,ms).
This directly addresses the concern that MAX-SAT is ``over-engineering'': the solver is essentially free; the cost lies in extraction.
At a marginal cost of 6.2 seconds per percentage point of accuracy gain, the overhead is modest relative to the improvement.


\subsection{Qualitative Analysis: Constraint Quality}
\label{sec:casestudies}

We examine 8 representative problems from the Qwen2.5-14B evaluation to assess constraint quality.
Of the extracted constraint sets, 85\% are semantically meaningful (correctly capturing logical relationships from the reasoning trace), and 75\% of constraints excluded by the MAX-SAT solver correspond to genuine reasoning errors.

We observe four case categories:
\begin{itemize}[nosep,leftmargin=*]
  \item \textbf{Correct constraints, correct answer} (3/8): Both SC and \sica{} succeed; constraints are well-formed and consistent.
  \item \textbf{Correct constraints, \sica{} rescues SC error} (2/8): SC ties between correct and incorrect answers; \sica{}'s constraint weighting correctly identifies the answer with higher global consistency.
  \item \textbf{Noisy constraints, \sica{} robust} (2/8): Some constraints are incorrect, but MAX-SAT's weighted optimization assigns them low weight, and the correct answer still wins.
  \item \textbf{Bad constraints, \sica{} fails} (1/8): Ambiguous negation scope (``not all'' $\to$ $\forall x\,\neg P(x)$ instead of $\neg\forall x\, P(x)$) leads to an incorrect MAX-SAT selection.
\end{itemize}

The failure mode is informative: negation scope ambiguity is a known weakness of FOL extraction from natural language, and transparent failure modes support future targeted improvement.


\subsection{Cross-Family Validation}
\label{sec:crossfamily}

To rule out the possibility that \sica{}'s gains are specific to the Qwen model family, we evaluate on Mistral-7B-Instruct-v0.3 and LLaMA-3.1-8B-Instruct (Table~\ref{tab:main}).

Both models show $\czero{} \approx 0.40$ and \sica{} gains of +5.0\,pp, with FOLIO improvements of +10.0\,pp each.
Tiebreak win rates (79\% for Mistral, 77\% for LLaMA) are comparable to Qwen-14B's 78\%.
In the $\czero{}$-vs-gap scatter plot, Mistral-7B ($\czero{}{=}0.42$, $\Delta{=}+5.0$) and LLaMA-8B ($\czero{}{=}0.40$, $\Delta{=}+5.0$) cluster together despite being independent model families---same capacity level, similar $\czero{}$, similar gain.
This is direct evidence that $\czero{}$ is a universal predictor rather than a family-specific artifact.


% ============================================================
\section{Discussion}
\label{sec:discussion}

\paragraph{$\czero{}$ and epistemic uncertainty.}
While we define $\czero{}$ operationally as multi-candidate rate, it is conceptually related to epistemic uncertainty: high $\czero{}$ indicates that the model's knowledge is insufficient to consistently determine a single answer.
We deliberately avoid formal uncertainty quantification terminology, as our measure captures sampling diversity rather than calibrated confidence.
The connection to epistemic uncertainty is suggestive but requires formal validation against established UQ benchmarks.

\paragraph{Practical deployment guidance.}
Our results suggest a simple deployment protocol: run SC once, compute $\czero{}$; if $\czero{}>20\%$, deploy \sica{} for an expected gain proportional to $\czero{}$.
For reasoning-specialized models or tasks where the model is highly confident ($\czero{}<5\%$), SC suffices.
This makes \sica{} a \emph{conditional} improvement---practitioners need not commit to the 38\% overhead without evidence that it will help.

\paragraph{Limitations.}
Our evaluation is limited to logic and math reasoning domains.
While the $\czero{}$ framework is domain-agnostic, constraint extraction quality may vary across domains---the GSM8K pilot suggests that numeric constraints are less expressive than FOL predicates.
The D012 subset evaluations (60 questions) for cross-family and scaling experiments lack statistical power for individual significance; their value is in the aggregate pattern across the $\czero{}$-gap scatter.
The constraint extraction step relies on the same model, inheriting its biases.

\paragraph{Why not just use a reasoning model?}
Our reasoning model results (QwQ-32B, DeepSeek-R1) show that stronger models don't need \sica{} because they are already self-consistent.
However, instruction-tuned models remain the dominant deployment paradigm due to lower cost, lower latency, and broader availability.
\sica{} specifically targets this deployed model class, providing gains precisely where they are most needed.


% ============================================================
\section{Conclusion}
\label{sec:conclusion}

We introduced \sica{}, a symbolic consensus aggregation method that extracts FOL constraints from LLM reasoning traces and uses weighted partial MAX-SAT to resolve voting ties.
\sica{} achieves 77\% tiebreak accuracy (vs.\ 50\% for random), yielding +6.9\,pp on FOLIO with Qwen2.5-14B ($p{=}0.023$).
We showed that \sica{}'s effectiveness is governed by trace diversity $\czero{}$, which predicts the accuracy gap with $r{=}0.91$ across 14 configurations spanning 3 model families, 2 reasoning architectures, and 2 domains.
The $\czero{}{=}20\%$ threshold cleanly separates the active regime (instruction models, challenging tasks) from the degenerate regime (reasoning models, easy tasks), enabling zero-cost deployment decisions.
MAX-SAT-based constraint aggregation outperforms simpler alternatives (+1.7\,pp over count-only) at negligible solver cost (1.4\% of wall time), confirming that satisfiability checking is essential, not over-engineering.


% ============================================================
\bibliographystyle{plain}

\begin{thebibliography}{10}

\bibitem{wang2023selfconsistency}
X.~Wang, J.~Wei, D.~Schuurmans, Q.~Le, E.~Chi, S.~Narang, A.~Chowdhery, and D.~Zhou.
\newblock Self-Consistency Improves Chain of Thought Reasoning in Language Models.
\newblock In \emph{ICLR}, 2023.

\bibitem{chen2023universal}
X.~Chen, Y.~Shi, J.~Wang, Z.~Li, and P.~Yin.
\newblock Universal Self-Consistency for Large Language Model Generation.
\newblock \emph{arXiv preprint arXiv:2311.17311}, 2023.

\bibitem{pan2023logiclm}
L.~Pan, A.~Albalak, X.~Wang, and W.~Wang.
\newblock Logic-LM: Empowering Large Language Models with Symbolic Solvers for Faithful Logical Reasoning.
\newblock In \emph{EMNLP Findings}, 2023.

\bibitem{olausson2023linc}
T.~X.~Olausson, A.~Gu, B.~Lipkin, C.~Zhang, A.~Solar-Lezama, J.~Tenenbaum, and R.~Levy.
\newblock LINC: A Neurosymbolic Approach for Logical Reasoning by Combining Language Models with First-Order Logic Provers.
\newblock In \emph{EMNLP}, 2023.

\bibitem{kuhn2023semantic}
L.~Kuhn, Y.~Gal, and S.~Farquhar.
\newblock Semantic Uncertainty: Linguistic Invariances for Uncertainty Estimation in Natural Language Generation.
\newblock In \emph{ICLR}, 2023.

\bibitem{xiong2024llms}
M.~Xiong, Z.~Hu, X.~Lu, Y.~Li, J.~Fu, J.~He, and B.~Hooi.
\newblock Can LLMs Express Their Uncertainty? An Empirical Evaluation of Confidence Elicitation in LLMs.
\newblock In \emph{ICLR}, 2024.

\bibitem{han2022folio}
S.~Han, H.~Wu, X.~Liu, J.~Pan, L.~Huang, and M.~Jiang.
\newblock FOLIO: Natural Language Reasoning with First-Order Logic.
\newblock \emph{arXiv preprint arXiv:2209.00840}, 2022.

\bibitem{tafjord2021proofwriter}
O.~Tafjord, B.~Dalvi, and P.~Clark.
\newblock ProofWriter: Generating Implications, Proofs, and Abductive Statements over Natural Language.
\newblock In \emph{ACL Findings}, 2021.

\bibitem{cobbe2021gsm8k}
K.~Cobbe, V.~Kosaraju, M.~Bavarian, M.~Chen, H.~Jun, L.~Kaiser, M.~Plappert, J.~Tworek, J.~Hilton, R.~Nakano, C.~Hesse, and J.~Schulman.
\newblock Training Verifiers to Solve Math Word Problems.
\newblock \emph{arXiv preprint arXiv:2110.14168}, 2021.

\end{thebibliography}

\end{document}
